\documentclass{ximera}



% MATH -----------------------------------------------------------
\newcommand{\nats}{\mathbb N}
\newcommand{\ints}{\mathbb Z}
\newcommand{\rats}{\mathbb Q}
\newcommand{\reals}{\mathbb R}
\newcommand{\complex}{\mathbb C}
\newcommand{\powerset}{\mathscr P}

%-------------------------------------------------------------

\pgfplotsset{compat=1.18}
\begin{document}
\begin{abstract}
 \end{abstract}

    \title{Exemplar Examples 22}
    
    \maketitle

\begin{enumerate}

\item Let's use linearity and $\mathcal{L}[y\,^{\prime\prime}]=s^2\mathcal{L}[y]-sy(0)-y\,^{\prime}(0)$ to derive a formula for $\mathcal{L}[\sin{(\omega t)}]$. 


\vfill


 \item Now let's use above and $\mathcal{L}[y\,^{\prime}]=s\mathcal{L}[y]-y(0)$ to derive a formula for $\mathcal{L}[\cos{(\omega t)}]$.


\vfill

\newpage


\item Let's use $\displaystyle \mathcal{L}(tf(t))=-\frac{dF}{ds}$ to derive formulas for $\mathcal{L}[t\cos{(\omega t)}]$ and $\mathcal{L}[t\sin{(\omega t)}]$.



\newpage

\item Let's solve the IVP $y\,^{\prime\prime}+y=4\cos{(t)}$, $y(0)=1$, $y\,^{\prime}(0)=0$ using Laplace transforms.

\begin{enumerate}
\item First let's find the Laplace transform $Y(s)$ of the solution $y(t)$.

\vfill

\item  Now let's  find the inverse Laplace transform of $Y(s)$, which is the solution $y(t)$.


\vfill
\vfill

% y= cos(t) + 2t sin(t).

% -------------------------------------------------------------------

% Laplace:

% s^2 Y - s +Y= 4s/(s^2+1)

% (s^2+1) Y =s +  4s/(s^2+1)

%  Y =s/(s^2+1) +  2 (2s)/(s^2+1)^2

%  y= cos(t) + 2 t sin(t) and DONE!!




\end{enumerate}


\end{enumerate}

\end{document}